\documentclass[envcountsect]{svmono}

% choose options for [] as required from the list
% in the Reference Guide

%\usepackage{mathptmx}
%\usepackage{helvet}
%\usepackage{courier}
%


\usepackage{type1cm}

\usepackage{makeidx}         % allows index generation
\usepackage{graphicx}        % standard LaTeX graphics tool
                             % when including figure files
\usepackage{multicol}        % used for the two-column index
\usepackage[bottom]{footmisc}% places footnotes at page bottom

\usepackage{newtxtext}       %
\usepackage{newtxmath}       % selects Times Roman as basic font

\usepackage[unicode]{hyperref}
\usepackage[nameinlink,capitalize]{cleveref}
\crefname{appendix}{Appendix}{Appendices}
\Crefname{appendix}{Appendix}{Appendices}
\crefname{problem}{Problem}{Problems}
\Crefname{problem}{Problem}{Problems}
\crefname{conjecture}{Conjecture}{Conjectures}
\Crefname{conjecture}{Conjecture}{Conjectures}
\hypersetup{
colorlinks=true,
linktoc=all
}
\makeatletter
\renewcommand*{\theHtheorem}{\theHsection.\arabic{theorem}}
\makeatother

\usepackage{enumitem}
\usepackage{mathtools}
\providecommand{\xlongrightarrow}[2][]{\xrightarrow[#1]{#2}}
\usepackage{tikz-cd}
\usetikzlibrary{decorations.pathreplacing,arrows.meta,calc,patterns}
\usepackage{dsfont}

\usepackage{fancyhdr}
\renewcommand{\headrulewidth}{0pt}

\usepackage[final]{microtype}
\usepackage{inputenc}
\usepackage{textcomp}
\usepackage[T1]{fontenc}

%\usepackage[color]{showkeys}
%\definecolor{refkey}{rgb}{0,0,1}


% see the list of further useful packages
% in the Reference Guide

\makeindex             % used for the subject index
                       % please use the style svind.ist with
                       % your makeindex program

%%%%%%%%%%%%%%%%%%%%%%%%%%%%%%%%%%%%%%%%%%%%%%%%%%%%%%%%%%%%%%%%%%%%%

\graphicspath{{./}{images/}}

\newcommand{\nef}{\mathrm{nef}}
\newcommand{\Bir}{\mathrm{Bir}}
\newcommand{\NS}{\mathrm{NS}}
\newcommand{\lex}{\mathrm{lex}}
\newcommand{\NA}{\mathrm{NA}}
\newcommand{\SH}{\mathrm{SH}}
\newcommand{\PSH}{\mathrm{PSH}}
\newcommand{\QPSH}{\mathrm{QPSH}}
\newcommand{\BM}{\mathrm{BM}}
\newcommand{\Conv}{\mathrm{Conv}}
\newcommand{\FS}{\mathrm{FS}}
\newcommand{\An}{\mathrm{An}}
\newcommand{\val}{\mathrm{val}}
\newcommand{\Macron}{\mathrm{Defect}}
\newcommand{\Bl}{\mathrm{Bl}}
\newcommand{\Ele}{\mathrm{Ele}}
\newcommand{\triv}{\mathrm{triv}}
\newcommand{\gf}{\mathrm{gf}}
\newcommand{\tor}{\mathrm{tor}}
\newcommand{\prop}{\mathrm{prop}}
\newcommand{\Star}{\mathrm{Star}}
\newcommand{\Haus}{\mathrm{Haus}}
\newcommand{\sg}{\mathrm{sg}}
\newcommand{\TC}{\mathrm{TC}}
\newcommand{\Kah}{\textup{Käh}}
\newcommand{\sups}{\operatorname*{sup*}}
\newcommand{\uscsup}[1][]{\operatorname*{sup}_{#1}{}\!{}^{*}}
\newcommand{\ddc}{\mathrm{dd}^{\mathrm{c}}}
\newcommand{\ddt}{\frac{\mathrm{d}}{\mathrm{d}t}}
\newcommand{\ddtz}[1][0]{\left.\ddt\right|_{t={#1}}}
\newcommand{\loc}{\mathrm{loc}}
\newcommand{\convC}{\xrightarrow{\mathrm{C}}}
\newcommand{\qsup}{\operatorname*{q-sup}}
\newcommand{\qinf}{\operatorname*{q-inf}}

\DeclareMathOperator{\Img}{Im}
\DeclareMathOperator{\Nm}{Norm}
\DeclareMathOperator{\nn}{nn}
\DeclareMathOperator{\nK}{nK}
\DeclareMathOperator{\Capa}{Cap}
\DeclareMathOperator{\Real}{Re}
\DeclareMathOperator{\Spec}{Spec}
\DeclareMathOperator{\Sing}{Sing}
\DeclareMathOperator{\Sgn}{Sign}
\DeclareMathOperator{\Rest}{Res}
\DeclareMathOperator{\DHm}{DH}
\DeclareMathOperator{\Fin}{Fin}
\DeclareMathOperator{\Reg}{Reg}
\DeclareMathOperator{\Tr}{Tr}
\DeclareMathOperator{\Trop}{Trop}
\DeclareMathOperator{\Div}{div}
\DeclareMathOperator{\ord}{ord}
\DeclareMathOperator{\Dom}{Dom}
\DeclareMathOperator{\epi}{epi}
\DeclareMathOperator{\LCE}{CE}
\DeclareMathOperator{\Supp}{Supp}
\DeclareMathOperator{\MA}{MA}
\DeclareMathOperator{\vol}{vol}
\DeclareMathOperator{\Pic}{Pic}
\DeclareMathOperator{\Int}{Int}
\DeclareMathOperator{\RelInt}{RelInt}
\DeclareMathOperator{\cl}{cl}
\DeclareMathOperator{\cor}{cor}
\DeclareMathOperator{\GL}{GL}
\DeclareMathOperator{\Aut}{Aut}
\DeclareMathOperator{\usc}{usc}
\DeclareMathOperator{\mult}{mult}
\DeclareMathOperator{\Bs}{Bs}




\setlist[enumerate,1]{label=\textup{(\arabic*)},ref=\textup{(\arabic*)}}

\makeatletter
% For labels of the form itm:<parent label>:<item>, print "Parent (item)".
\newcommand{\itemref}[1]{\itemref@split\cref#1:\@nil}
\newcommand{\Itemref}[1]{\itemref@split\Cref#1:\@nil}
\def\itemref@split#1itm:#2:#3:#4:\@nil{#1{#2:#3}\,\ref{itm:#2:#3:#4}}
\makeatother



% --- Hyphenation hints for technical terms (added during indexing pass) ---
\hyphenation{
  plu-ri-sub-har-mon-ic plu-ri-sub-har-mon-ic-i-ty plu-ri-po-lar plu-ri-po-ten-tial
  pseu-do-con-vex pseu-do-met-ric pseu-do-norm sub-har-mon-ic
  bi-mero-mor-phic mero-mor-phic ho-lo-mor-phic
  semi-con-tin-u-ous semi-con-ti-nu-i-ty up-per-semi-con-ti-nu-i-ty
  multi-plier multi-pli-er ana-lyt-i-fi-ca-tion
  equi-sin-gu-lar equi-sin-gu-lar-i-ty quasi-equi-sin-gu-lar
  trop-i-cal-i-za-tion non-Arch-i-me-de-an Arch-i-me-de-an
  sub-ge-o-des-ic sub-ge-o-des-ics ge-o-des-ic ge-o-des-ics
  Mon-ge--Am-pe-re Käh-ler Lel-ong Bo-uck-som Ohsa-wa--Take-go-shi
  Da-r-vas Bed-ford--Tay-lor Don-ald-son Ross--Witt-Nys-tröm
  Bouck-som--Eys-sid-ieux--Guedj--Zer-i-ahi
  ho-mo-ge-ne-ous ho-mo-ge-ne-i-ty Berk-o-vich
  Brel-ot-Grau-ert--Rem-mert Hi-ron-a-ka Wein-stock
  Lazars-feld Ka-veh--Khov-an-skii Kis-el-man Choq-uet
}
% --- end hyphenation hints ---

%\includeonly{author/chapter_pshfunctions,author/chapter_npp,author/chapter_envelopes,author/chapter_Bergman}
%\includeonly{author/chapter_envelopes,author/chapter_comp,author/chapter_Bergman}



\usepackage{needspace}
\newcommand{\dc}{\mathrm d^{\mathrm c}}
\newcommand{\R}{\mathbb R}
\newcommand{\C}{\mathbb C}
\hypersetup{pdftitle={Nonuniqueness for a quadratically vanishing density},pdfauthor={},pdfsubject={AI-assisted mathematical working draft}}
\begin{document}
\author{}
\title{Nonuniqueness for a quadratically\\vanishing density}
\subtitle{Smooth solutions on a nonclosed background}
\date{Working draft, \today}
\pagenumbering{Alph}
\maketitle
\mainmatter
\chapter{Smooth solutions across a degenerate hypersurface}

This note collects constructions and proofs from the working manuscript \emph{Singularities in Hermitian pluripotential theory}. Using three rank-one directions, it gives a smooth nonclosed Hermitian background on $\mathbb P^3$ and a smooth density which vanishes exactly quadratically on two real hypersurfaces, is positive elsewhere, and has uncountably many distinct normalized smooth solutions.

The original construction is reproduced first so that the note can be read without the rest of the book. A first extension retains its background and removes the open zero-density band, with infinite-order vanishing on the remaining hypersurfaces. A second extension changes the background coefficients to achieve quadratic vanishing. Both arguments use a scalar ordinary differential equation, two endpoint barriers and estimates at every derivative order.

Let $X=\mathbb P^3$, with its Fubini--Study form $\omega_X$. For a real function $u$ on $X$, we use $\ddc=\frac{\mathrm{i}}{2\pi}\partial\bar\partial$ and $\dc=\frac{\mathrm{i}}{4\pi}(\bar\partial-\partial)$. For a smooth real $(1,1)$-form $\omega$, a function is $\omega$-plurisubharmonic when $\omega+\ddc u\geq0$ as a current. For the smooth potentials constructed here, $\mathrm{MA}_\omega(u)=(\omega+\ddc u)^3$ is an ordinary smooth top form.

\begin{definition}\label{sec:obstacle-envelopes}
The canonical minimal potential is
\[
 V_\omega=\left(\sup\{v\in\PSH(X,\omega):v\leq0\}\right)^*.
\]
A potential has minimal singularities if its difference from $V_\omega$ is bounded. When $\omega>0$, we have $V_\omega=0$, so every bounded potential has minimal singularities.
\end{definition}

For two smooth potentials $u,v$, put $A=\omega+\ddc u$ and $B=\omega+\ddc v$. The difference of their measures satisfies the algebraic identity
\begin{equation}\label{uni-eq:linearization}
 A^3-B^3=\ddc(u-v)\wedge(A^2+A\wedge B+B^2).
\end{equation}
The constructions below arrange equality of the cubes while leaving freedom in the potentials where the common density vanishes.

\begin{definition}\label{def:bmp}
A Hermitian form $\omega$ has the bounded mass property if the masses $\int_X(\omega+\ddc u)^3$ of bounded $\omega$-plurisubharmonic potentials have a common finite upper bound.
\end{definition}

The bounded-mass proposition below uses bounded Bedford--Taylor products and closed Stokes.

\begin{proposition}\label{prop:bmp-metric}
Every Hermitian form $\omega$ on $\mathbb P^3$ has bounded mass: there is $M_\omega<\infty$ such that
\[
 \int_X(\omega+\ddc u)^3\leq M_\omega
 \quad\text{for every bounded }u\in\PSH(X,\omega).
\]
\end{proposition}
\begin{proof}
Choose $C>0$ with $\omega\leq C\omega_X$. Then $C\omega_X+\ddc u\geq\omega+\ddc u\geq0$. Positivity of bounded mixed products and closed Stokes give
\[
 \int_X(\omega+\ddc u)^3
 \leq\int_X(C\omega_X+\ddc u)^3
 =C^3\int_X\omega_X^3.
\]
\end{proof}

\Needspace{12\baselineskip}
\section{Smooth nonuniqueness for a degenerate density}
\label{sec:degenerate-hermitian-nonuniqueness}

The first example puts the freedom in the potential inside an open region where the density is zero. Three independent rank-one directions reduce equality of the Monge--Amp\`ere measures to equality of scalar products. We construct two semipositive forms with the same cube, then add an exact correction to obtain a strictly positive equation background.

Let $X=\mathbb P(\mathbb C^2\oplus\mathbb C^2)=\mathbb P^3$, with block homogeneous coordinates $[A:B]$, where $A,B\in\mathbb C^2$, and let $\omega_X$ be its Fubini--Study form.

\begin{theorem}\label{thm:degenerate-hermitian-nonuniqueness}
On $X=\mathbb P^3$ there exist a smooth Hermitian form $\omega$ with $\mathrm{d}\omega\ne0$, a smooth nonnegative nonzero function $f$, and distinct smooth functions $u,v\in\PSH(X,\omega)$ such that
\[
 \sup_Xu=\sup_Xv=0,\qquad
 (\omega+\ddc u)^3=f\omega_X^3=(\omega+\ddc v)^3.
\]
The density vanishes on a nonempty open set. Both potentials have minimal singularities.
\end{theorem}

\subsection{Three rank-one directions}
\label{subsec:numa-rank-directions}

On $U=\{A\ne0,B\ne0\}$ put
\[
 r=\log\frac{|A|^2}{|B|^2},\qquad
 P=\ddc\log|A|^2,\quad Q=\ddc\log|B|^2,\quad
 R=\mathrm{d}r\wedge \dc r.
\]
The logarithms individually use a local homogeneous lift; their Hessians are independent of that lift. The forms $P,Q$ are pullbacks of Fubini--Study forms from the two projective lines. All three forms are semipositive of rank one, and
\begin{equation}\label{eq:numa-rank-frame}
 P^2=Q^2=R^2=0,\qquad P\wedge Q\wedge R>0.
\end{equation}
Indeed, in local coordinates $A=(1,a)$ and $B=\tau(1,b)$, $P$ and $Q$ are positive multiples of $\mathrm{i}\,\mathrm{d}a\wedge \mathrm{d}\bar a$ and $\mathrm{i}\,\mathrm{d}b\wedge \mathrm{d}\bar b$, while
\[
 \partial r=\frac{\bar a\,\mathrm{d}a}{1+|a|^2}
             -\frac{\mathrm{d}\tau}{\tau}
             -\frac{\bar b\,\mathrm{d}b}{1+|b|^2}.
\]
\Needspace{6\baselineskip}
Its $\mathrm{d}\tau$ component is nonzero. The three rank-one directions are therefore independent. In particular, for nonnegative $x,y,z$,
\begin{equation}\label{eq:numa-cube}
 (xP+yQ+zR)^3=6xyz\,P\wedge Q\wedge R,
\end{equation}
and the form is positive exactly when all three coefficients are positive.

\Needspace{5\baselineskip}
For a smooth function $G:\mathbb R\to\mathbb R$ the chain rule is
\begin{equation}\label{eq:numa-chain}
 \ddc G(r)=G'(r)(P-Q)+G''(r)R.
\end{equation}
\Needspace{7\baselineskip}
Let $\sigma(r)=\mathrm{e}^r/(1+\mathrm{e}^r)$. Since
$\log(|A|^2+|B|^2)=\log|B|^2+\log(1+\mathrm{e}^r)$, we have
\begin{equation}\label{eq:numa-fs}
 \omega_X=\sigma(r)P+\sigma(-r)Q+\sigma'(r)R.
\end{equation}
A function of $r$ constant at each end of $\mathbb R$ extends smoothly over $X\setminus U$. A form whose coefficients equal those of \eqref{eq:numa-fs} for $|r|\geq2$ extends smoothly by $\omega_X$ there.

\subsection{The scalar coefficients}
\label{subsec:numa-scalar-coefficients}

The central coefficients will make a different rank-one factor vanish in each solution form. Outside the central interval they will return to the Fubini--Study coefficients, so the resulting forms extend across the omitted projective lines. Primes denote derivatives with respect to $r$.

Choose an even smooth bump $\beta$ supported on $[-1,1]$, strictly positive on $(-1,1)$ and with $\beta'>0$ on $(-1,0)$, normalized by $\int_{-1}^1\beta=2$. A positive multiple of $\exp(-1/(1-r^2))$ on $(-1,1)$, extended by zero, does this. Define
\[
 F(r)=-1+\int_{-1}^r\beta(t)\,\mathrm{d}t\quad(-1<r<1),
 \qquad F=-1\ (r\leq-1),\quad F=1\ (r\geq1),
\]
and put $p=\beta/2$, $q=\beta'/2$. Then $F$ is smooth and odd, $p$ is even, and $q$ is odd, positive on $(-1,0)$ and negative on $(0,1)$. Both $p$ and $q$ vanish to infinite order at the endpoints.

Take smooth $a\geq0$ with $a=1$ on $(-\infty,-2/3]$ and $a=0$ on $[-1/3,\infty)$, and put $b(r)=a(-r)$. Choose an even nonnegative $\chi\in C_c^\infty((-1/3,1/3))$ with $\chi(0)>0$. On $[-1,1]$ set
\[
 z_0=(q^2+\chi^2)^{1/2}.
\]
It is smooth: the expression under the square root is positive in the interior; near each endpoint $z_0=|q|$ with fixed sign of $q$. It is flat at the endpoints and positive in the interior. Also,
\begin{equation}\label{eq:numa-coefficient-cancellation}
 z_0\geq|q|,\qquad a(z_0-q)=0,\qquad b(z_0+q)=0.
\end{equation}
The cancellations follow from the disjoint supports of $a$ and $\chi$, and from the sign of $q$ on the support of $a$; the other equality follows by reflection.

\Needspace{12\baselineskip}
On $(-\infty,-1]$ choose smooth functions $\ell,s$ such that $0\leq\ell\leq1$, $\ell=0$ on $(-\infty,-2]$, $\ell=1$ on $[-3/2,-1]$, while $s>0$ for $r<-1$, $s=1$ on $(-\infty,-3/2]$, and $s$ is flat with value zero at $-1$. Such functions are obtained from ordinary smooth cutoffs and a flat positive endpoint function. Define
\[
 \begin{aligned}
 x&=(1-\ell)\sigma+\ell,&
 y&=(1-\ell)\sigma(-r)+\ell s,&
 z&=(1-\ell)\sigma'+\ell s &&(r\leq-1),\\
 x&=p+a,& y&=p+b,& z&=z_0 &&(-1\leq r\leq1),
 \end{aligned}
\]
and for $r\geq1$ use $x(r)=y(-r)$, $y(r)=x(-r)$, $z(r)=z(-r)$.

The definitions join smoothly. At $-1$ their values are $(1,0,0)$, with all discrepancies flat; near that endpoint $a=1$ and $b=0$. The central definitions are symmetric under reflection, so the same holds at $1$. We have
\[
 x,y,z>0\quad\hbox{for }r\ne\pm1,\qquad
 (x,y,z)(-1)=(1,0,0),\quad (x,y,z)(1)=(0,1,0).
\]
For $|r|\geq2$ their values are $(\sigma(r),\sigma(-r),\sigma'(r))$.

\subsection{Two forms with the same cube}
\label{subsec:numa-equal-cubes}

The cancellation identities now turn the scalar coefficients into two semipositive forms with the same cube. Their average will be strictly positive except at the two ends of the central interval.

The globally smooth form and function
\[
 \omega_0=x(r)P+y(r)Q+z(r)R,\qquad w=F(r)
\]
are understood by the extensions above. The form is semipositive, and positive except on $M_- =\{r=-1\}$ and $M_+=\{r=1\}$. Put
\[
 H=\tfrac12\ddc w=p(P-Q)+qR,\qquad A_0=\omega_0-H,\quad B_0=\omega_0+H.
\]
For $|r|\geq1$ we have $H=0$. On $|r|<1$,
\[
 A_0=aP+(2p+b)Q+(z_0-q)R,\qquad
 B_0=(2p+a)P+bQ+(z_0+q)R.
\]
All coefficients are nonnegative. The cancellations in \eqref{eq:numa-coefficient-cancellation} give $A_0^3=B_0^3=0$ on this band. Outside it both cubes equal $\omega_0^3$. Thus $A_0,B_0$ are globally smooth semipositive forms with equal cubes. That common top-degree form is smooth, vanishes on the closed band $|r|\leq1$, and is positive on its complement, including near $X\setminus U$.

\Needspace{12\baselineskip}
\subsection{Restoring strict positivity of the background}
\label{subsec:numa-positive-background}

The forms $A_0,B_0$ already have the required common measure. We obtain a Hermitian background by adding the same exact correction to their average and subtracting its potential from both solutions. The correction must restore positivity at $M_\pm$ while preserving it in the transition regions.

Take $\kappa\in C_c^\infty((-2,2))$, equal to one on $[-3/2,3/2]$. For a parameter $\varepsilon>0$ to be chosen, set
\[
 \rho(r)=\varepsilon(r^2-1)\kappa(r),\qquad
 \omega=\omega_0+\ddc\rho.
\]
Its coefficients on $U$ are $x+\rho'$, $y-\rho'$ and $z+\rho''$. On $[-3/2,3/2]$ one has $\rho'=2\varepsilon r$ and $\rho''=2\varepsilon$. The last coefficient is positive. For $r\geq0$, $x+2\varepsilon r$ is positive, including at $r=0$ since $x(0)=p(0)>0$. On $[-3/2,0]$, the function $x$ has a positive minimum $m$, so $x+2\varepsilon r\geq m-3\varepsilon>0$ for small $\varepsilon$. Reflection gives the corresponding estimate for $y-2\varepsilon r$. On the compact region $3/2\leq|r|\leq2$, all three original coefficients have a positive minimum and the perturbations are $O(\varepsilon)$. Outside that region and the central interval the form is exactly $\omega_X$. This proves strict positivity globally.

For completeness, nonclosedness can be checked directly. Since $\mathrm{d}R=-\mathrm{d}r\wedge(P-Q)$,
\begin{equation}\label{eq:numa-differential}
 \mathrm{d}\omega=\mathrm{d}\omega_0
   =\mathrm{d}r\wedge\bigl((x'-z)P+(y'+z)Q\bigr).
\end{equation}
At $r=0$, $x'=y'=0$ and $z=\chi(0)>0$. Hence $\mathrm{d}\omega=z\,\mathrm{d}r\wedge(Q-P)\ne0$ there. The nonzero $\mathrm{d}\tau$ component used in \eqref{eq:numa-rank-frame} also shows that $\mathrm{d}r\wedge P$ and $\mathrm{d}r\wedge Q$ are linearly independent.

\begin{proof}[of \cref{thm:degenerate-hermitian-nonuniqueness}]
Set $\widetilde u=-w/2-\rho$ and $\widetilde v=w/2-\rho$. The construction gives
\[
 \omega+\ddc\widetilde u=A_0\geq0,\qquad
 \omega+\ddc\widetilde v=B_0\geq0.
\]
Subtract their respective suprema to obtain $u,v$. Their Hessians are unchanged, and $v-u=w+\text{constant}$ is nonconstant. The function $f=A_0^3/\omega_X^3=B_0^3/\omega_X^3$ is smooth, nonnegative and nonzero. It vanishes on the nonempty open band $|r|<1$. Both smooth potentials have minimal singularities because $V_\omega=0$ for a Hermitian background.
\end{proof}

Every Hermitian background on $\mathbb P^3$ has bounded mass: the Fubini--Study form has invariant mass by closed bounded Stokes, and \cref{prop:bmp-metric} makes finiteness independent of the metric. This applies to all the projective backgrounds constructed below.

The linearized difference also explains the failure of the smooth maximum argument. The tensor
\[
 T=A_0^2+A_0\wedge B_0+B_0^2
\]
satisfies $\ddc w\wedge T=B_0^3-A_0^3=0$, but $T=0$ on $M_\pm$, where $A_0=B_0=\omega_0$ has rank one. The resulting equation is not globally elliptic, even though both potentials and their common density are smooth.

\section{Nonuniqueness with a zero set of empty interior}
\label{sec:numa-thin-zero-set}

We now keep the Hermitian background of \cref{thm:degenerate-hermitian-nonuniqueness}, but change the density so that it is positive off $M_-\cup M_+$, where $M_\pm=\{r=\pm1\}$. For potentials of the form $G(r)-\rho$, the equation on the middle band becomes a first-order ordinary differential equation for the slope $G'$. Barriers keep its solutions between the admissible slopes, and infinite-order vanishing at the endpoints then makes their extensions by zero smooth.

\Needspace{9\baselineskip}
\begin{theorem}\label{thm:numa-thin-zero-set}\index{Zero set!thin degeneracy set}
There are a smooth nonclosed Hermitian form $\omega$ on $\mathbb P^3$ and a smooth nonnegative density $g$, positive off two disjoint smooth real hypersurfaces, such that the equation
\[
 \mathrm{MA}_\omega(u)=g\omega_X^3
\]
has uncountably many distinct smooth solutions normalized by $\sup_Xu=0$. The background $\omega$ has bounded mass.
\end{theorem}

Thus a zero set with empty interior, even one of measure zero, does not ensure uniqueness of smooth normalized solutions.

\Needspace{12\baselineskip}
We use the coefficients of \cref{subsec:numa-scalar-coefficients}, with the explicit choice
\[
 p(r)=C\exp\!\left(-\frac1{1-r^2}\right),\qquad -1<r<1,
 \qquad \int_{-1}^1p(r)\,\mathrm{d}r=1,
\]
extended by zero, where the integral condition determines $C>0$; then $q=p'$. Fix $0<s<1$, and set
\[
 k_s(r)=(x(r)+sp(r))(y(r)-sp(r))(z(r)+sq(r)).
\]
For $|r|>1$ this is $xyz$. If $G'=sp$, then the potential $G(r)-\rho$ has $\omega+\ddc(G(r)-\rho)=\omega_0+sH$. This form is strictly positive for $-1<r<1$, since
\[
 x+sp=a+(1+s)p>0,\quad
 y-sp=b+(1-s)p>0,\quad
 z_0+sq\geq(1-s)z_0>0.
\]
It is also strictly positive outside the closed band, and at $r=\pm1$ it equals the rank-one form $\omega_0$. Consequently
\begin{equation}\label{eq:numa-thin-density}
 g\omega_X^3=(\omega_0+sH)^3
\end{equation}
defines a smooth density $g$ that is positive exactly off $M_-\cup M_+$. These hypersurfaces have measure zero and empty interior. On the band $|r|<1$, a potential $G(r)-\rho$ with slope $h=G'$ solves the equation exactly when $(x+h)(y-h)(z+h')=k_s$, which is the differential equation \eqref{eq:numa-thin-ode} below.

\Needspace{14\baselineskip}
\begin{lemma}\label{lem:numa-thin-trapping}
There exists $\eta_0>0$ such that for every $0<\eta<\eta_0$ the initial-value problem
\begin{equation}\label{eq:numa-thin-ode}
 h'=\frac{k_s(r)}{(x(r)+h)(y(r)-h)}-z(r),
 \qquad h(0)=sp(0)-\eta,
\end{equation}
has a smooth solution on $(-1,1)$ with $-x<h<y$ and $h<sp$. Moreover, on some fixed endpoint neighborhoods there are constants $0<c_L<s<t<1$, independent of $\eta$, such that
\[
 c_Lp<h<sp\quad\hbox{near }-1,
 \qquad -tp<h<sp\quad\hbox{near }1.
\]
\end{lemma}
\begin{proof}
The function $h_s=sp$ solves \eqref{eq:numa-thin-ode} and lies strictly between $-x$ and $y$ throughout $(-1,1)$. Write $F_s(r,h)$ for the right-hand side of the equation, which is smooth on the open domain $-x(r)<h<y(r)$.

Near $-1$ the coefficients are exactly
\[
 x=1+p,\qquad y=p,\qquad z=q>0.
\]
For a constant $c\in(-1,1)$, the product of coefficients for the trial slope $cp$ is
\[
 k_c=(1-c^2)\bigl(1+(1+c)p\bigr)pq.
\]
Fix $c_L\in(0,s)$. At $p=0$, the factor preceding $pq$ for $c=c_L$ exceeds that for $c=s$ by $s^2-c_L^2>0$. Thus, on a sufficiently small fixed neighborhood of the left endpoint, $k_{c_L}>k_s$ and
\begin{equation}\label{eq:numa-left-trap}
 F_s(r,c_Lp)-(c_Lp)'=
 \frac{k_s-k_{c_L}}{(x+c_Lp)(y-c_Lp)}<0.
\end{equation}
Thus $c_Lp$ is a lower barrier when the equation is solved backward, that is, toward $-1$.

Near $1$ the exact coefficients are
\[
 x=p,\qquad y=1+p,\qquad z=-q>0,
\]
and the trial slope $cp$ has product of coefficients
\[
 k_c=(1-c^2)\bigl(1+(1-c)p\bigr)p(-q).
\]
Fix $s<t<1$. At $p=0$, the factor for $c=s$ exceeds that for $c=-t$ by $t^2-s^2>0$. On a sufficiently small fixed neighborhood of the right endpoint, therefore, $k_s>k_{-t}$, so
\begin{equation}\label{eq:numa-right-trap}
 F_s(r,-tp)-(-tp)'=
 \frac{k_s-k_{-t}}{(x-tp)(y+tp)}>0.
\end{equation}
Thus $-tp$ is a lower barrier when the equation is solved forward, toward $1$.

Choose $r_L<0<r_R$ in these respective neighborhoods. The graph of $h_s$ over $[r_L,r_R]$ is a compact subset of the domain of $F_s$. Existence and continuous dependence on initial values for smooth ordinary differential equations therefore give $\eta_0>0$ such that, for every $0<\eta<\eta_0$, the solution with initial value $sp(0)-\eta$ exists on $[r_L,r_R]$, stays in the domain, and satisfies
\[
 h(r_L)>c_Lp(r_L),\qquad h(r_R)>-tp(r_R).
\]
By local uniqueness, two solutions of the equation cannot cross. Since $h(0)<h_s(0)$, we have $h<h_s$ wherever $h$ exists.

Follow the solution backwards from $r_L$; it cannot cross $c_Lp$. Indeed, at a first crossing point the difference $h-c_Lp$ would vanish and be positive immediately to its right, so its derivative there would be nonnegative, whereas \eqref{eq:numa-left-trap} shows that this derivative is strictly negative. Thus $c_Lp<h<sp$ on the part of the interval of existence to the left of $r_L$. Likewise, \eqref{eq:numa-right-trap} prevents a first downward crossing of $-tp$ to the right of $r_R$, giving $-tp<h<sp$ there.

These bounds also show that the maximal interval of existence is all of $(-1,1)$. Indeed, near a possible endpoint inside $(-1,1)$, the factors $x+h$ and $y-h$ stay bounded away from zero, so the graph of $h$ stays in a compact subset of the domain of $F_s$, and the continuation theorem for ordinary differential equations extends the solution beyond that endpoint. Hence the solution exists on $(-1,1)$ and satisfies the claimed bounds.
\end{proof}

\Needspace{7\baselineskip}
\begin{lemma}\label{lem:numa-thin-flat-extension}
Every solution furnished by \cref{lem:numa-thin-trapping}, extended by zero outside $(-1,1)$, is smooth on $\mathbb R$ and vanishes to infinite order at $\pm1$.
\end{lemma}
\begin{proof}
Let $h$ be one of the solutions in \cref{lem:numa-thin-trapping}. On either endpoint neighborhood put $\xi=h/p$; since $p>0$ on $(-1,1)$, this quotient is smooth there, and the trapping bounds place $\xi$ in a fixed compact subinterval of $(-1,1)$. Define the rational functions
\[
 R_L(p,\xi)=
 \frac{(1-s^2)(1+(1+s)p)}{(1+(1+\xi)p)(1-\xi)}-1,
\]
\[
 R_R(p,\xi)=1-
 \frac{(1-s^2)(1+(1-s)p)}{(1+\xi)(1+(1-\xi)p)}.
\]
The explicit formulas for the coefficients near the endpoints give
\[
 h'=q\,R_L(p,\xi)\quad\hbox{on the left},\qquad
 h'=q\,R_R(p,\xi)\quad\hbox{on the right}.
\]
Both rational functions are smooth, with all derivatives bounded, on a neighborhood of the relevant set of pairs $(p,\xi)$ with $p\geq0$ small and $\xi$ in the compact interval above. Writing $R_{L/R}$ for the function corresponding to the endpoint under consideration, we obtain from $h'=p'\xi+p\xi'$ and $q=p'$
\begin{equation}\label{eq:numa-scaled-ode}
 \xi'=\frac{p'}p\bigl(R_{L/R}(p,\xi)-\xi\bigr).
\end{equation}
For the chosen bump,
\[
 \frac{p'}p=-\frac{2r}{(1-r^2)^2}.
\]
This rational function and each of its derivatives grow at most like a power of the inverse distance to the endpoint. Moreover, each derivative of $p$ is $p$ times a rational function with the same property.

An induction using \eqref{eq:numa-scaled-ode} shows that every derivative $\xi^{(j)}$ grows at most polynomially in the inverse distance to the endpoint. The case $j=0$ is the trapping bound on $\xi$. For the induction step, differentiate the right-hand side of \eqref{eq:numa-scaled-ode}, including its term $-\xi$, $j$ times. Every resulting term is a product of a derivative of $p'/p$, a bounded partial derivative of $R_{L/R}$, and derivatives of $p$ and of $\xi$ of order at most $j$. All of these factors have polynomial bounds, by the induction hypothesis for the derivatives of $\xi$ and since every derivative of $p$ is bounded. This gives the bound for $\xi^{(j+1)}$.

Leibniz' rule for $h=p\xi$ now gives, for each $j$, constants $C_j,N_j$ such that
\[
 |h^{(j)}(r)|\leq C_j p(r)(1-|r|)^{-N_j}
\]
near the endpoints. The right-hand side tends to zero at the endpoints, because $p$ decays exponentially in the inverse distance to the endpoint. Thus all derivatives of $h$ tend to zero at $\pm1$. Extend each derivative by zero outside $(-1,1)$. By the mean value theorem, the difference quotient of $h^{(j)}$ at an endpoint tends to the limit of $h^{(j+1)}$, which is zero. Hence these extensions are the successive derivatives of the extension of $h$, which is therefore smooth and vanishes to infinite order at $\pm1$.
\end{proof}

\begin{proof}[of \cref{thm:numa-thin-zero-set}]
Let $\eta_0$ be as in \cref{lem:numa-thin-trapping}, and for $0<\eta<\eta_0$ denote by $h_\eta$ the solution with initial value $sp(0)-\eta$. Extend $h_\eta$ by zero, which gives a smooth function by the preceding lemma, and take a primitive $G_\eta$. It is constant near each end of the real line, so $G_\eta(r)$ extends smoothly to all of $\mathbb P^3$, including the two omitted projective lines. Set
\[
 \widetilde u_\eta=G_\eta(r)-\rho.
\]
On the open band,
\[
 \omega+\ddc\widetilde u_\eta
 =(x+h_\eta)P+(y-h_\eta)Q+(z+h_\eta')R.
\]
The coefficients of $P$ and $Q$ are strictly positive by \cref{lem:numa-thin-trapping}, while the equation gives
\[
 z+h_\eta'=\frac{k_s}{(x+h_\eta)(y-h_\eta)}>0.
\]
Outside the band the form is $\omega_0$, and on the two boundary hypersurfaces it is also the rank-one form $\omega_0$, since $h_\eta$ vanishes to infinite order there. It is therefore semipositive on $X$, and by \eqref{eq:numa-cube} its cube equals $6k_sP\wedge Q\wedge R=(\omega_0+sH)^3$ on the band and $\omega_0^3=(\omega_0+sH)^3$ elsewhere. Thus its cube is everywhere the measure $g\omega_X^3$ of \eqref{eq:numa-thin-density}. Finally, normalize by subtracting $\sup_X\widetilde u_\eta$.

Distinct values of $\eta$ give different slopes at $r=0$, so the corresponding potentials do not differ by a constant, and their normalized versions are distinct. This gives an uncountable family. The background is that of \cref{thm:degenerate-hermitian-nonuniqueness}, so it has bounded mass.
\end{proof}

\Needspace{8\baselineskip}
The density vanishes to infinite order on the two hypersurfaces. Their complement has three components, and on each of the two outer components, two different solutions differ by a constant, with different constants on the two components. The differential equation on the middle component joins these constants smoothly across the degenerate hypersurfaces. This is why ellipticity at every point of the positive-density set does not give a maximum principle argument across its boundary.

\section{Quadratic density zeros}
\label{sec:quadratic-density-nonuniqueness}

Infinite-order vanishing of the density is not needed for smooth nonuniqueness. We change the coefficients of $\omega_0$ so that its coefficient of $R$ stays positive and one of its coefficients of $P$ and $Q$ vanishes quadratically at each of the two hypersurfaces $M_\pm$. The resulting singular differential equation forces the slopes of differences of solutions to decay like $\exp(-A/d)$, where $d$ is the distance to the hypersurface and $A>0$. The slopes therefore extend smoothly by zero, even though the density vanishes only to finite order.

\Needspace{14\baselineskip}
\begin{theorem}\label{thm:quadratic-density-nonuniqueness}\index{Zero set!quadratic vanishing}
On $\mathbb P^3$ there exist a smooth nonclosed Hermitian form $\omega$ and a smooth nonnegative density $g$ such that:
\begin{enumerate}
\item The zero set of $g$ consists of two disjoint smooth real hypersurfaces. Near either hypersurface, $g$ is the square of a smooth defining function times a smooth strictly positive function.
\item The equation $(\omega+\ddc u)^3=g\omega_X^3$ has uncountably many distinct smooth $\omega$-plurisubharmonic solutions normalized by $\sup_Xu=0$.
\end{enumerate}
The background $\omega$ has bounded mass.
\end{theorem}

We keep the rank-one forms $P,Q,R$ and the function $r$ of \cref{subsec:numa-rank-directions}, and use new scalar coefficients.

Choose a smooth function $a$ valued in $[0,1]$, equal to one for $r\leq-2/3$ and zero for $r\geq-1/3$, and put $b(r)=a(-r)$. Choose a smooth even cutoff $\chi$ valued in $[0,1]$, zero for $|r|\leq5/4$ and one for $|r|\geq3/2$. Set $p(r)=(1-r^2)^2$ and
\begin{equation}\label{eq:quadratic-background}
 \begin{split}
 x&=(1-\chi)(p+a)+\chi\sigma,\\
 y&=(1-\chi)(p+b)+\chi(1-\sigma),\\
 z&=1-\chi+\chi\sigma',\qquad
 \omega_0=xP+yQ+zR.
 \end{split}
\end{equation}
Since these coefficients coincide with those of \eqref{eq:numa-fs} for $|r|\geq3/2$, the form $\omega_0$ extends by $\omega_X$ across the two omitted projective lines. The coefficient $z$ is strictly positive, and $x,y$ are strictly positive except that $y(-1)=0$ and $x(1)=0$. Thus $\omega_0$ is positive off $M_\pm=\{r=\pm1\}$ and has rank two there.

Choose a smooth cutoff $\kappa\in C_c^\infty((-2,2))$ equal to one on $[-3/2,3/2]$. For a parameter $\varepsilon>0$ to be chosen, define the correction
\[
 \rho=\varepsilon(r^2-1)\kappa(r),\qquad
 \omega=\omega_0+\ddc\rho.
\]
For sufficiently small $\varepsilon>0$, the form $\omega$ is strictly positive. Indeed, on $[-3/2,3/2]$ its coefficients are $x+2\varepsilon r$, $y-2\varepsilon r$ and $z+2\varepsilon$. The first is positive for $r\geq0$, and on $[-3/2,0]$ the function $x$ has a positive minimum, so it remains positive for small $\varepsilon$. Reflection gives the positivity of the second coefficient, and that of the third is immediate. On $3/2\leq|r|\leq2$, the form $\omega_0$ is strictly positive on a compact set and the correction is $O(\varepsilon)$. Near the projective ends there is no correction. The form $\omega$ is not closed: at $r=0$, $x'=y'=0$ and $z=1$, so by \eqref{eq:numa-differential}
\[
 \mathrm{d}\omega=\mathrm{d}r\wedge\bigl((x'-z)P+(y'+z)Q\bigr)\ne0.
\]

Define the density $g$ by
\[
 g\omega_X^3=\omega_0^3.
\]
It is smooth and positive exactly off $M_-\cup M_+$. By \eqref{eq:numa-fs} and \eqref{eq:numa-cube}, $g=xyz/(\sigma(r)(1-\sigma(r))\sigma'(r))$. Near either hypersurface, $xyz=(1+p)p$, and hence
\[
 g(r)=\frac{(1+p(r))(1-r^2)^2}{\sigma(r)(1-\sigma(r))\sigma'(r)}.
\]
The denominator is smooth and strictly positive at $\pm1$, and $1-r^2$ is a defining function for $M_\pm$ near these hypersurfaces. This proves the asserted quadratic vanishing. The potential $-\rho$ is one smooth solution.

\Needspace{15\baselineskip}
\begin{lemma}\label{lem:quadratic-density-trapping}
For every sufficiently small $\eta>0$, the equation
\begin{equation}\label{eq:quadratic-density-ode}
 h'=\frac{xyz}{(x+h)(y-h)}-z,\qquad h(0)=-\eta
\end{equation}
has a smooth solution on $(-1,1)$ with $-x<h<y$ and $h<0$. There are a constant $0<c<1$ and fixed neighborhoods of both endpoints, independent of $\eta$, on which
\[
 -cp<h<0.
\]
Every such solution, extended by zero, is a smooth function on $\mathbb R$ that vanishes to infinite order at $\pm1$.
\end{lemma}
\begin{proof}
The right-hand side $F(r,h)$ is smooth on $-x<h<y$, and $h=0$ is a solution. Near $-1$, $(x,y,z)=(1+p,p,1)$; near $1$, $(x,y,z)=(p,1+p,1)$. Fix $0<c<1$. For the trial function $-cp$, near the left endpoint
\[
 F(r,-cp)+cp'
 =\frac{1+p}{(1+(1-c)p)(1+c)}-1+cp'
 \longrightarrow-\frac{c}{1+c}<0,
\]
whereas near the right endpoint
\[
 F(r,-cp)+cp'
 =\frac{1+p}{(1-c)(1+(1+c)p)}-1+cp'
 \longrightarrow\frac{c}{1-c}>0.
\]
as $r\to\pm1$, since $p,p'\to0$ at both endpoints. The strict inequalities therefore hold on fixed neighborhoods of the endpoints, and $-cp$ is a lower barrier for the equation solved backwards near $-1$ and forwards near $1$.

Choose a compact interval $[r_L,r_R]\ni0$ with endpoints in those neighborhoods. Continuous dependence on initial values, around the solution $h=0$, gives $\eta_0>0$ such that for every $0<\eta<\eta_0$ the solution exists on $[r_L,r_R]$, stays in the domain of $F$, and lies above $-cp$ at $r_L$ and $r_R$. By uniqueness of solutions, it cannot meet the solution zero, so $h<0$. The barrier argument used in the proof of \cref{lem:numa-thin-trapping} then gives $-cp<h<0$ near both ends. On each compact subinterval of $(-1,1)$ the trapped graph stays at a positive distance from the boundary of the domain of $F$; hence, by the continuation theorem, the solution exists on all of $(-1,1)$.

\Needspace{6\baselineskip}
The equation factors as
\begin{equation}\label{eq:quadratic-density-logarithm}
 h'=h\,L(r,h),\qquad
 L(r,h)=\frac{z(x-y+h)}{(x+h)(y-h)}.
\end{equation}
Since $h<0$, we have $(\log(-h))'=L(r,h)$. Near the left endpoint, the trapping bounds and the explicit coefficients imply
\[
 L(r,h)=\frac{1+h}{(1+p+h)(p-h)}\geq\frac{C_1}{p}
\]
after shrinking that neighborhood, with $C_1>0$. Near the right endpoint they give
\[
 L(r,h)=\frac{-1+h}{(p+h)(1+p-h)}\leq-\frac{C_2}{p},
 \qquad C_2>0.
\]
Indeed, on the left, $1+h\geq1-cp\geq1/2$ and the denominator is at most $(1+p)(1+c)p$; on the right, the numerator is at most $-1$, and the denominator is at most $p(1+(1+c)p)$. Since $p$ is comparable to the square of the distance $d=1-|r|$ to the endpoint, integrating these inequalities backwards on the left and forwards on the right gives, for some $A>0$,
\begin{equation}\label{eq:quadratic-density-exponential}
 |h(r)|\leq C \mathrm{e}^{-A/d}
\end{equation}
near either endpoint. The constants may depend on the solution; uniformity in $\eta$ is not needed for the smooth extension.

It remains to control all derivatives. On the trapped graph, all partial derivatives in $(r,h)$ of the rational function $L$ grow at most polynomially in $d^{-1}$. Indeed, near an endpoint $x,y$ are polynomials in $r$, $z=1$, $h$ is bounded, and of the two factors in the denominator, one is bounded below by a positive constant and the other by $(1-c)p$. Repeated differentiation of the quotient therefore only introduces finite powers of $p^{-1}$.

Inductively, for every $j\geq0$ there are constants $C_j,N_j$ with
\[
 |h^{(j)}(r)|\leq C_j \mathrm{e}^{-A/d}d^{-N_j}.
\]
The case $j=0$ is \eqref{eq:quadratic-density-exponential}. Differentiate $h'=hL(r,h(r))$ $j$ times. Every resulting term is the product of a derivative of the factor $h$, of order at most $j$, and a derivative of the composite function $L(r,h(r))$. The latter involves only partial derivatives of $L$ and derivatives of $h$ of order at most $j$, which are polynomially bounded by the induction hypothesis, while the first factor has the exponential bound. This proves the case $j+1$. All derivatives of $h$ therefore tend to zero at $\pm1$. Extending them by zero and applying the mean value theorem to the difference quotients at the endpoints, as in \cref{lem:numa-thin-flat-extension}, shows that the extension of $h$ by zero is smooth and vanishes to infinite order at $\pm1$.
\end{proof}

\begin{proof}[of \cref{thm:quadratic-density-nonuniqueness}]
Let $\eta_0>0$ be such that \cref{lem:quadratic-density-trapping} applies for $0<\eta<\eta_0$, and let $h_\eta$ denote the solution with $h_\eta(0)=-\eta$. Take a primitive $G_\eta$ of its extension by zero. Since $G_\eta$ is constant near each end of the real line, $G_\eta(r)$ defines a smooth function on $X$. The potential $u_\eta=G_\eta(r)-\rho$ satisfies
\[
 \omega+\ddc u_\eta
 =(x+h_\eta)P+(y-h_\eta)Q+(z+h_\eta')R.
\]
Inside the band, the first two coefficients are positive by the lemma, and the equation gives
\[
 z+h_\eta'=\frac{xyz}{(x+h_\eta)(y-h_\eta)}>0.
\]
Outside the band the form is $\omega_0$, and on the two boundary hypersurfaces it is also the rank-two form $\omega_0$, since $h_\eta$ vanishes to infinite order there. Thus it is semipositive on $X$, and by \eqref{eq:numa-cube} its cube is $6xyzP\wedge Q\wedge R=\omega_0^3=g\omega_X^3$. Different initial values give potentials that do not differ by constants, so after normalization by the supremum we obtain an uncountable family. Bounded mass follows from \cref{prop:bmp-metric}, as for \cref{thm:degenerate-hermitian-nonuniqueness}.
\end{proof}

At each of the hypersurfaces $M_\pm$, the difference of two of these solutions agrees to infinite order with a constant, because its derivative is a difference of slopes that vanish to infinite order there. The constants at the two hypersurfaces need not agree. Thus the singular differential equation produces differences that are flat at $M_\pm$, even though the density has only quadratic zeros.


\section{Connected positive-density sets}
The positive-density sets of both hypersurface constructions above have three components. By contrast, the book proves uniqueness for bounded solutions with the same smooth nonnegative density and supremum zero when its positive set is connected and dense, on any connected compact Hermitian manifold. Its proof uses the energy of the potential difference and smooth test potentials concentrated near an extremal level. The topology of the positive-density region therefore separates these examples from that uniqueness theorem.

\section{Historical note}
The stability theorem \cite[Theorem 4.1]{KN21} assumes a uniform lower bound $f\geq c_0>0$ in addition to its capacity hypotheses. The densities in \cref{thm:numa-thin-zero-set} have no such lower bound. The stability theorem therefore does not apply to these examples.

\backmatter
\bibliographystyle{sgptalpha}
\bibliography{references}
\end{document}
